% TOP_PROBLEM: {"id": 1497, "title": "Infinitude of Perfect Numbers", "category_id": 1, "category": {"id": 1, "name": "number_theory", "display_name": "Number Theory", "description": "Properties of integers, prime numbers, Diophantine equations.", "slug": "number-theory", "order_index": 1, "created_at": "2026-07-31T15:26:25.670Z"}, "status": "open"}
% FIRST_POSTED: null
% ENTRY_KIND: statement_only
% Imported from: batch08.tex; UnsolvedMath ID 1497
\documentclass[11pt]{article}
\usepackage[margin=25mm]{geometry}
\usepackage{fontspec}
\setmainfont{Latin Modern Roman}
\usepackage{amsmath,amssymb,mathtools}
\usepackage{unicode-math}
\setmathfont{Latin Modern Math}
\usepackage{xurl,hyperref}
\hypersetup{colorlinks=true,urlcolor=blue}
\setcounter{secnumdepth}{0}
\setlength{\parindent}{0pt}
\setlength{\parskip}{5pt}
\setlength{\emergencystretch}{3em}
\newcommand{\sref}[2]{\hyperlink{source:#1:#2}{[S#2]}}
\newcommand{\cataloguescope}[1]{\paragraph{Recorded target.}#1}
\begin{document}
% UnsolvedMath ID 1497; source catalogue code NUM-003
\setcounter{section}{1496}
\section{Infinitude of Perfect Numbers}\label{1497}
{\small\sffamily UnsolvedMath ID 1497 · Number Theory}
\subsection{Definitions and mathematical statement}

\paragraph{Shared notation.}$\mathbb N=\{1,2,\ldots\}$, and $\mathbb Z,\mathbb Q,\mathbb R,\mathbb C$ denote the integers, rationals, reals and complex numbers. Any other conventions are specified locally.


For \(n\in\mathbb N\), let
\[
 D(n)=\{d\in\mathbb N:d\mid n\},\qquad
 \sigma(n)=\sum_{d\in D(n)}d.
\]
The proper positive divisors are \(D(n)\setminus\{n\}\).
A perfect number satisfies \(\sum_{d\mid n,\,d<n}d=n\), equivalently
\(\sigma(n)=2n\).
The Euclid--Euler theorem classifies even perfect numbers as
\(2^{p-1}(2^p-1)\), where \(p\) and \(2^p-1\) are prime.
This gives an equivalence between infinitely many even perfect numbers
and infinitely many Mersenne primes. The present question also allows
odd perfect numbers.
\paragraph{Statement recorded in the supplied source.}
Is the set
\[
                    \{n\in\mathbb N:\sigma(n)=2n\}
\]
infinite? Equivalently, for every \(B\in\mathbb N\), does there exist
a perfect number \(n>B\)? \sref{1497}{2}
\subsection{Short English statement}
Are there infinitely many perfect numbers, including both possible parities?
\subsection{Sources}
\begin{description}
\item[\hypertarget{source:1497:1}{S1}] ULAM.AI and the UnsolvedMath Contributors, \emph{UnsolvedMath: A Curated Collection of Open Mathematics Problems}, record 1497 (NUM-003). CC BY 4.0. \url{https://huggingface.co/datasets/ulamai/UnsolvedMath}.
\item[\hypertarget{source:1497:2}{S2}] John Cowles and Ruben Gamboa, \emph{Perfect Numbers in ACL2}, EPTCS 192 (2015), 53--59, DOI 10.4204/EPTCS.192.5; arXiv:1509.06081v1, Sections 1--2, the unrestricted infinitude question and $\sigma(n)=2n$ definition. \url{https://arxiv.org/abs/1509.06081v1}.
\end{description}
\label{end:1497}
\end{document}
